% SWQ: independent derivation, 2026-09-23.
\section{The original adelic trace quotient with every support label retained}
\label{sec:cc-supported-weil-quotient}

\paragraph{Human sources, objects, and reading scope.}
The analytic input is Alain Connes, Caterina Consani, and Matilde
Marcolli, \href{https://arxiv.org/abs/math/0703392v1}{\emph{The Weil proof
and the geometry of the adeles class space}}, original author file
\texttt{Yuri70v2.tex}, SHA-256
\texttt{b1f694cb934dd4fef829287f1a55dd8de21c67a481e552453c2f7e8fe0ae6fb5},
canonical programme source \texttt{PUBUNIT-5DE2A2D44F8B1470AC701FDA}.
The sections actually read for this derivation are source lines
610--710, 1316--1425, 1508--1551, 1680--1830, 1980--2160,
2161--2216, and 2313--2424, together with the title and author lines.
The original locators are \texttt{SCK}, \texttt{rangeTrpi},
\texttt{rangecVdef}, \texttt{varthetaH1V},
\texttt{varthetafconvol}, \texttt{vanishV},
\texttt{tracepairing}, \texttt{fsharpinv},
\texttt{TraceformulaThm}, and \texttt{TraceformulaH1}.
The paragraph following \texttt{motH1AKCK} specifies closure of
the range in the cokernel. The trace-class and explicit-formula
theorems are cited analytic inputs. The convolution ideal, its sharp
invariance with all endpoint terms, and the full-label quotient and
receiving maps are derived below. No positivity assertion or RH
conclusion is made.

\subsection{The original test algebra and its domains}

Retain the source's unscaled test algebra and operations:
\begin{equation}\tag{SWQ1}
\begin{gathered}
 C=\mathbb A_{\mathbb Q}^{\times}/\mathbb Q^{\times},\qquad
 B=\bigcap_{\beta\in\mathbb R}|\cdot|^\beta\mathcal S(C),\\
 (f*g)(x)=\int_C f(h)g(h^{-1}x)\,d^{\times}h,\qquad
 f^{\sharp}(x)=|x|^{-1}\overline{f(x^{-1})}.
\end{gathered}
\end{equation}
For $\mathbb Q$, the actual continuous multiplicative section in
the author's subsection \texttt{XKCKsect} is
\begin{equation}\tag{SWQ2}
 C=\mathbb R_{>0}\times K,\qquad K=\widehat{\mathbb Z}^{\times},
 \qquad (r,u)\longmapsto(r,(u_p)_p),
 \qquad |(r,u)|=r.
\end{equation}
Indeed one multiplies an idele by the unique rational number that
makes every finite component a unit and its real component positive.
The ambiguity would have all finite valuations zero and positive real
component, so would be the rational number $1$. Multiplication
respects these representatives. The multiplicative Haar measure is
$dr/r\,du$, where $\int_Kdu=1$. The additive adelic measure is real
Lebesgue measure times the measures of volume $1$ on each
$\mathbb Z_p$. These exact measures are retained throughout.

The topology used for closure below is the usual weighted Bruhat
test-function topology: at each fixed open subgroup $H\subset K$,
functions are $H$-invariant and have all weighted Schwartz seminorms;
one then uses the Bruhat inductive-limit topology over these finite
levels. An explicit defining family at a fixed level, with $r=e^x$, is
\begin{equation}\tag{SWQ3}
 p_{N,m,j}(f)=\sup_{x\in\mathbb R,\ u\in K}
 e^{N|x|}(1+|x|)^m
 \left|\partial_x^j f(e^x,u)\right|,
 \quad N,m,j\in\mathbb Z_{\ge0}.
\end{equation}
These seminorms retain all real weights in SWQ1: weights $e^{\beta x}$
are bounded by $e^{N|x|}$ when $N\ge|\beta|$, and conversely the two
weights $e^{Nx}$ and $e^{-Nx}$ control $e^{N|x|}$; Leibniz's formula
controls their derivatives. This is a description of the original
test domain, not a change of function or scale in its operations.

All convolution integrals in SWQ1 converge absolutely. Using
$x=y+(x-y)$ and $1+|x|\le(1+|y|)(1+|x-y|)$ gives the explicit bound
\begin{equation}\tag{SWQ4}
 p_{N,m,j}(f*g)
 \le 2p_{N+1,m,0}(f)\,p_{N,m,j}(g).
\end{equation}
The constant $2$ is $\int_{\mathbb R}e^{-|y|}dy$; the compact Haar
factor is $1$. Differentiation can pass under the integral by the
same bound on every derivative. Convolution preserves a common
finite level, and this proves the continuity needed below. Since
$C$ is abelian, substitution $h\mapsto xh^{-1}$ proves commutativity.
Fubini, justified by the corresponding three-function $L^1$ bound,
proves associativity.

In these coordinates sharp is exactly
\begin{equation}\tag{SWQ5}
 f^{\sharp}(e^x,u)=e^{-x}\overline{f(e^{-x},u^{-1})},\qquad
 p_{N,m,j}(f^{\sharp})\le
 \sum_{\ell=0}^j\binom j\ell p_{N+1,m,\ell}(f).
\end{equation}
The bound follows by differentiating $e^{-x}$ and $f(e^{-x},u^{-1})$;
the two derivative signs together have absolute value $1$.
Direct substitution gives $(f^{\sharp})^{\sharp}=f$.
Inversion preserves Haar measure on this abelian group, and the norm
is a character. Substitution in the absolutely convergent convolution
therefore gives $(f*g)^{\sharp}=g^{\sharp}*f^{\sharp}$.
In particular sharp is continuous on the stated test space.
The algebra $B$ has no convolution unit among its functions.
Indeed let $b(e^x,u)=e^{-x^2}$, which belongs to $B$ and has
trivial compact-unit character. Its real Fourier transform is
$\sqrt\pi e^{-t^2/4}$, nonzero for every real $t$. A convolution
unit $a\in B$ would consequently have
$\int_{\mathbb R\times K}a(e^x,u)e^{itx}\,dx\,du=1$ for every
$t$. Integration by parts in $x$, justified by SWQ3, makes this
integral tend to zero as $|t|\to\infty$, a contradiction.
No unit distribution is inserted in this chapter.

\subsection{The adelic source, with both Poisson endpoints}

Use the canonical additive character trivial on $\mathbb Q$, with
real component $e^{-2\pi ixy}$ and finite conductor $\mathbb Z_p$,
and its self-dual measures specified after SWQ2. Write
\begin{equation}\tag{SWQ6}
\begin{gathered}
 \widehat\xi(y)=\int_{\mathbb A}\xi(x)\alpha(xy)\,dx,
 \qquad \varepsilon_0(\xi)=\xi(0),\qquad
 \varepsilon_1(\xi)=\int_{\mathbb A}\xi(x)\,dx,\\
 \mathcal S(\mathbb A)_0=\ker\varepsilon_0\cap\ker\varepsilon_1,
 \qquad (E\xi)(g)=\sum_{q\in\mathbb Q^{\times}}\xi(qg),
 \qquad V=E\mathcal S(\mathbb A)_0.
\end{gathered}
\end{equation}
For the reference characters $e_v$ in SWQ28, choose exactly
these local components of $\alpha$. This fixes more than their
conductors: it specifies the characters themselves, so that
the differential idele there is $a_v=1$ at every place.
The sum is invariant under the rational choice of representative.
For fixed $g$, finite-place compact support puts its contributing
rational $q$ in $D^{-1}\mathbb Z$ for one positive integer $D$.
Real Schwartz decay then makes the sum absolutely convergent,
locally uniformly with every real derivative. Unit dilation at
finite places preserves the required common denominator and bounds.

For completeness the full Poisson formula used here is
\begin{equation}\tag{SWQ7}
 \sum_{q\in\mathbb Q}\xi(qg)
   =|g|^{-1}\sum_{q\in\mathbb Q}\widehat\xi(qg^{-1}).
\end{equation}
It follows directly for the original Bruhat space. A finite-place
test function is a finite linear combination of characteristic
functions of cosets $a+N\widehat{\mathbb Z}$, with rational $a$ and
positive integer $N$, after a common denominator is chosen. The
rational elements in such a coset are $a+N\mathbb Z$. Real
periodization of a Schwartz function gives
\[
 \sum_{n\in\mathbb Z}\phi(a+Nn)
  =\frac1N\sum_{m\in\mathbb Z}
       e^{2\pi i ma/N}\widehat\phi(m/N).
\]
To verify this last identity, periodize over $N\mathbb Z$, integrate
over one period to compute every Fourier coefficient, and evaluate
the resulting absolutely convergent Fourier series at $a$.
The finite Fourier transform of the characteristic coset is
$N^{-1}\alpha_f(ay)1_{N^{-1}\widehat{\mathbb Z}}(y)$.
At rational $y=m/N$, its character is $e^{2\pi i am/N}$ because
the global character is trivial on rational numbers. These exact
factors prove adelic Poisson by linearity. Finally replacing the
test function by $x\mapsto\xi(gx)$ multiplies its Fourier transform
by $|g|^{-1}$ and evaluates it at $g^{-1}y$, proving SWQ7.
The same real Fourier inversion and finite character orthogonality
give $\widehat{\widehat\xi}(x)=\xi(-x)$.

Separating the actual $q=0$ terms in SWQ7, with neither endpoint
discarded, gives
\begin{equation}\tag{SWQ8}
 E\xi(g)=|g|^{-1}E\widehat\xi(g^{-1})
             +|g|^{-1}\varepsilon_1(\xi)-\varepsilon_0(\xi).
\end{equation}
For $|g|\ge1$, the earlier common-denominator estimate gives
$E\xi(g)=O_M(|g|^{-M})$ for every $M$, uniformly in the compact
unit coordinate and with every logarithmic real derivative.
For $\xi\in\mathcal S(\mathbb A)_0$, both displayed endpoints
vanish by definition, and the right side applies the same large-norm
estimate to $\widehat\xi$ at $g^{-1}$. Indeed
$\widehat\xi(0)=\varepsilon_1(\xi)=0$ and
$\int\widehat\xi=\xi(0)=0$ by Fourier inversion.
The factor $|g|^{-1}$ stays present. Since $M$ is arbitrary, it
proves all the small-norm bounds in SWQ3 as well. Thus $V\subset B$
on the source's actual domain.

Apply SWQ8 to $\widehat{\overline\xi}$, or equivalently apply
SWQ7 to $\overline\xi$ and invert $g$. The exact result is
\begin{equation}\tag{SWQ9}
 (E\xi)^{\sharp}(g)
 =E\widehat{\overline\xi}(g)
      -|g|^{-1}\overline{\varepsilon_0(\xi)}
      +\overline{\varepsilon_1(\xi)}.
\end{equation}
This identity retains the conjugations, the factor $|g|^{-1}$,
and both endpoints. Put $J\xi=\widehat{\overline\xi}$. Its
endpoints are
\begin{equation}\tag{SWQ10}
 \varepsilon_0(J\xi)=\overline{\varepsilon_1(\xi)},\qquad
 \varepsilon_1(J\xi)=\overline{\varepsilon_0(\xi)}.
\end{equation}
Consequently $J$ preserves the actual space $\mathcal S(\mathbb A)_0$
and SWQ9 gives $V^{\sharp}\subset V$; applying sharp twice gives
equality. This derives sharp invariance from additive Fourier--Poisson,
rather than assuming that the original quotient was an involutive
quotient.

\subsection{The actual convolution ideal and its closure}

For $f\in B$, $\xi\in\mathcal S(\mathbb A)_0$, use the actual
multiplicative section SWQ2 to define
\begin{equation}\tag{SWQ11}
 \eta(x)=\int_C f(h)\xi(h^{-1}x)\,d^{\times}h.
\end{equation}
This integral belongs to the original Bruhat--Schwartz space.
Indeed the orbit of the finite-place component of $\xi$ under
the compact group $K$ lies in a finite-dimensional space of test
functions: its compact support stays in one finite product of
bounded balls, and its common local-constancy subgroups are unchanged
by units. The finite quotient of this compact support by these
subgroups has finitely many points. For the real component, every
Schwartz seminorm has the explicit bound
\begin{equation}\tag{SWQ12}
 \sup_{x_\infty,x_f}
 |x_\infty^N\partial_{x_\infty}^j\eta(x)|
 \le C_{\xi,N,j}\int_C |f(h)|\,|h|^{N-j}\,d^{\times}h<\infty.
\end{equation}
The last integral is finite by SWQ3. Differentiation and integration
are justified by the same bound, while the fixed finite-place
test space makes the resulting function a Bruhat test function.

The two endpoints of this integral are exactly
\begin{equation}\tag{SWQ13}
 \eta(0)=\varepsilon_0(\xi)\int_C f(h)\,d^{\times}h,\qquad
 \int_{\mathbb A}\eta(x)\,dx
 =\varepsilon_1(\xi)\int_C |h|f(h)\,d^{\times}h.
\end{equation}
For the second equation the absolute double integral is bounded by
$\|\xi\|_{L^1(\mathbb A)}\int_C|h||f(h)|d^{\times}h$, which is
finite. The change of variables $x=hy$ therefore gives its exact
adelic factor $|h|$, and Fubini is valid for this equation.
Thus $\eta\in\mathcal S(\mathbb A)_0$.

There is also a justified interchange with the arithmetic sum in
the convolution calculation. For fixed $g$, real Schwartz decay and
the fixed rational denominator give, uniformly in the unit part of
$h$, the estimate
\[
 \sum_{q\ne0}|\xi(qh^{-1}g)|\le C_{\xi,g}(1+|h|).
\]
For example, dominate the real values on the grid $D^{-1}\mathbb Z$
by $C(1+|n||g|/(D|h|))^{-2}$, and bound the resulting series by
its first term plus its integral. Multiplying by $|f(h)|$ is
integrable by SWQ3. Hence
\begin{equation}\tag{SWQ14}
 f*E\xi=E\eta\in V.
\end{equation}
Commutativity gives the ideal property in both orders. This proves
the actual ideal assertion needed for the quotient, including the
source-domain check in Lemma \texttt{varthetaH1V}.

Let
\begin{equation}\tag{SWQ15}
 W=\overline V\subset B,\qquad Q=B/W,
\end{equation}
where closure is in the topology stated before SWQ3. Continuity of
convolution by each fixed test function, established in SWQ4,
makes $W$ a convolution ideal. Continuity and involutivity of
sharp in SWQ5, together with SWQ9--10, give $W^{\sharp}=W$.
Thus
$[f]*[g]=[f*g]$ and $[f]^{\sharp}=[f^{\sharp}]$
are well-defined on $Q$. Notice that SWQ13 is not an exchange of
$\sum_{q\ne0}$ with the degree integral of $E\xi$. It does not
assert $\int_C|g|E\xi(g)d^{\times}g=0$. The author's
\texttt{fubini1}--\texttt{nothappen} calculation concerns precisely
that different, generally invalid exchange.

\subsection{The full-label quotient and the different Bourne quotient}

Let $L$ be any bounded distributive lattice. All constructions below
hold without finiteness; to require $e\ne\tau$ take $L$ nontrivial.
Define the original constrained carrier and operations
\begin{equation}\tag{SWQ16}
\begin{gathered}
 G_L(B)=\{(0,\lambda):\lambda\in L\}\cup(B\times\{1_L\}),\\
 (f,\lambda)+(g,\mu)=(f+g,\lambda\vee\mu),\qquad
 (f,\lambda)*(g,\mu)=(f*g,\lambda\wedge\mu),\\
 (f,\lambda)^{\sharp}=(f^{\sharp},\lambda),\qquad
 \tau=(0,0_L),\quad e=(0,1_L).
\end{gathered}
\end{equation}
This is a commutative nonunital convolution semiring with sharp.
Carrier closure follows because any nonzero output amplitude requires
all relevant nonzero input amplitudes to have top label. For the
product this uses that convolution with a zero function is zero;
it does not assume that $B$ has no zero divisors. Associativity and
distributivity combine those of convolution with the lattice laws.
Sharp is an involution preserving labels and reversing convolution.
The complex scalar action has the separate semimodule formula
$(c,\alpha)(f,\lambda)=(cf,\alpha\wedge\lambda)$ for
$(c,\alpha)\in G_L(\mathbb C)$.

The exact same-label quotient is
\begin{equation}\tag{SWQ17}
 (f,\lambda)\sim_W(g,\mu)\quad\Longleftrightarrow\quad
 \lambda=\mu\ \hbox{and}\ f-g\in W,
 \qquad G_L(B)/\!\sim_W\ \simeq G_L(Q).
\end{equation}
To prove this, use the map
$\pi_L(f,\lambda)=([f],\lambda)$. It respects addition,
convolution, sharp, and the scalar action by the ideal and
involution properties just proved. It is onto: a nonzero quotient
amplitude has top label and lifts to a representative with top label;
a zero amplitude with any label lifts to $(0,\lambda)$.
Its fibres are exactly the displayed relation. Alternatively, adding
a third element preserves equality of labels and preserves the
amplitude difference modulo $W$; multiplication sends that difference
to an element of $W$; sharp sends it into $W$ as well. These direct
checks prove the required congruence and the quotient assertion.

In particular every $(0,\lambda)$ survives, and $e$ and $\tau$
remain distinct for nontrivial $L$. A nonzero $w\in W$ with top
label maps to $e$, which differs from $\tau$ when $L$ is nontrivial.
For nontrivial $L$, the inverse image of the absorbing
additive zero $\tau$ under $\pi_L$ is only $\{\tau\}$, even though
$\pi_L$ can identify different top-labelled amplitudes. Thus the
preimage of a single zero does not specify this semiring congruence.
For the one-element lattice instead $G_L(B)=B$, and that inverse
image is the ordinary amplitude kernel $W$.

The ordinary Bourne congruence by the ideal $I=G_L(W)$ is different:
\begin{equation}\tag{SWQ18}
 x\equiv_I y\quad\Longleftrightarrow\quad
 \exists a,b\in I:\ x+a=y+b,
 \qquad
 (f,\lambda)\equiv_I(g,\mu)\Longleftrightarrow f-g\in W.
\end{equation}
The forward implication follows from the amplitude equation.
Conversely, if $f-g\in W$, choose
$a=(g-f,1_L)$ and $b=(0,1_L)$. Both belong to $I$; the two
amplitudes become $g$ and both labels become $1_L$, proving equality.
Consequently
\begin{equation}\tag{SWQ19}
 G_L(B)/\!\equiv_I\ \simeq Q,
 \qquad
 G_L(B)/\!\sim_W\ \longrightarrow Q,
 \quad([f],\lambda)\longmapsto[f].
\end{equation}
The Bourne quotient collapses the entire labelled-zero fibre, so it
identifies $e$ and $\tau$. More exactly it is the same-label quotient
followed by the congruence generated by $e=\tau$: multiplication
of this equality by $(0,\lambda)$ makes $(0,\lambda)=\tau$ for
every $\lambda$. After these identifications the carrier and
operations are exactly those of $Q$. No Boolean character or
selected support point has been used.

\subsection{Every numerical additive trace and the full-labelled receiver}

Let $A$ be any abelian group. Every additive map
$t:G_L(B)\to A$ kills every labelled zero: its idempotence gives
$t(0,\lambda)+t(0,\lambda)=t(0,\lambda)$, hence $t(0,\lambda)=0$.
Define $\ell(f)=t(f,1_L)$. Then $\ell$ is additive on $B$, since
addition of two top-labelled elements stays top-labelled and
$\ell(0)=t(e)=0$. Every nonzero amplitude already has top label,
so
\begin{equation}\tag{SWQ20}
 t(f,\lambda)=\ell(f),\qquad
 \operatorname{Hom}_{+}(G_L(B),A)
       \simeq\operatorname{Hom}_{+}(B,A).
\end{equation}
Conversely composition of any additive $\ell$ with amplitude
projection is additive, proving both directions. This includes all
numerical additive traces into $\mathbb C$, not only continuous or
complex-linear ones. Cyclicity adds no restriction here because
convolution is commutative. Such a map descends through SWQ17
exactly when $\ell(W)=0$: necessity follows by comparing $(w,1_L)$
with $e$, and sufficiency follows from the displayed difference
condition. The same classification on $G_L(Q)$ is
$\operatorname{Hom}_{+}(G_L(Q),A)=\operatorname{Hom}_{+}(Q,A)$.
Thus numerical additivity itself removes zero-label distinctions;
it is not a proof that those distinctions were absent from the
source object.

Retain the original numerical cohomological trace
\begin{equation}\tag{SWQ21}
 T(f)=\operatorname{Tr}\bigl(\underline\vartheta_m(f)\mid H^1\bigr).
\end{equation}
Here $\underline\vartheta_m(f)$ is the original convolution action
on the source's quotient realization, with its prescribed closure.
Its trace-class meaning is the cited Theorem
\texttt{varthetagTrH1}, with the trace formula
\texttt{TraceformulaThm}; this chapter does not replace that
analytic theorem by a finite-dimensional trace assumption.
Its vanishing on $W$ can be checked on the actual quotient action:
for $w\in W$ and $b\in B$, $w*b\in W$, hence
$\underline\vartheta_m(w)[b]=0$. The operator is zero, so
\begin{equation}\tag{SWQ22}
 T(W)=0,\qquad \overline T:Q\to\mathbb C,\quad
 \overline T([f])=T(f).
\end{equation}
This does not require an unproved exchange of an infinite spectral
sum with a limit of representatives: the operator itself vanishes.

The original trace and Weil-pairing receivers retaining the labels are
\begin{equation}\tag{SWQ23}
\begin{gathered}
 T_L:G_L(Q)\longrightarrow G_L(\mathbb C),\qquad
 T_L([f],\lambda)=(T(f),\lambda),\\
 \mathcal W_L\bigl(([f],\lambda),([g],\mu)\bigr)
   =\bigl(T(f*g^{\sharp}),\lambda\wedge\mu\bigr).
\end{gathered}
\end{equation}
The first map is well-defined by SWQ22. It respects addition and
the declared $G_L(\mathbb C)$ scalar action; it is a trace receiver,
not a claimed multiplicative character. For the second, replace
$f$ by $f+w$ and $g$ by $g+v$, with $w,v\in W$. Its change of
amplitude is
\begin{equation}\tag{SWQ24}
 T\bigl(w*g^{\sharp}+f*v^{\sharp}+w*v^{\sharp}\bigr)=0.
\end{equation}
Every summand belongs to $W$ by SWQ15. This proves representative
independence while retaining each term of the actual change.
If one input quotient amplitude is zero, the output amplitude is
zero. Thus every nonzero output amplitude has top label and the
receiver belongs to the declared constrained carrier.

The receiver is additive in both variables and complex-linear in
the first, conjugate-linear in the second, with the exact lattice
factors. For example, if $a=(c,\alpha)$, then
\begin{equation}\tag{SWQ25}
 \mathcal W_L(ax,y)=a\mathcal W_L(x,y),\qquad
 \mathcal W_L(x,ay)=\overline a\,\mathcal W_L(x,y),
 \quad\overline a=(\overline c,\alpha).
\end{equation}
The amplitudes follow from convolution and sharp. For additivity,
the label equality is
$(\lambda\vee\lambda')\wedge\mu
=(\lambda\wedge\mu)\vee(\lambda'\wedge\mu)$, and likewise in
the second variable. These prove the assertions directly.
In particular
$\mathcal W_L((0,\lambda),([g],\mu))=(0,\lambda\wedge\mu)$.
Even a zero numerical amplitude therefore retains its exact support
in SWQ23. Projecting to $\mathbb C$ forgets that support by SWQ20.

\subsection{The complete original trace terms of this receiver}

For the trivial unit character write
$M f(s)=\int_C f(u)|u|^s d^{\times}u$. Every such integral
converges absolutely on the full complex plane by SWQ3, locally
uniformly with every derivative. Haar inversion gives
\begin{equation}\tag{SWQ26}
 M(f^{\sharp})(s)=\overline{M f(1-\overline s)},\qquad
 M(f*g^{\sharp})(s)=M f(s)\overline{M g(1-\overline s)}.
\end{equation}
For the second identity the weighted double integral is finite,
so multiplicative Fubini is valid. At the identity one has instead
the exact weighted inner-product expression
\begin{equation}\tag{SWQ27}
 (f*g^{\sharp})(1)
    =\int_C f(u)\overline{g(u)}|u|\,d^{\times}u.
\end{equation}
In particular $M(f^{\sharp})(1)=\overline{M f(0)}$.
The conjugation is necessary for complex tests; it is missing in
the author's displayed \texttt{codegcorrZf} equality if that equality
is read for unrestricted complex-valued $f$. No real-valued
restriction is imposed here.

Put $h=f*g^{\sharp}$. The source's full trace formula therefore
gives the following amplitude for SWQ23:
\begin{equation}\tag{SWQ28}
\begin{split}
 T(h)={}&M f(0)\overline{M g(1)}
       +M f(1)\overline{M g(0)}\\
 &-(\log|a|)\int_C f(u)\overline{g(u)}|u|d^{\times}u\\
 &-\sum_{v\le\infty}\int_{(\mathbb Q_v^{\times},e_v)}'
                \frac{h(u^{-1})}{|1-u|_v}\,d_v^{\times}u.
\end{split}
\end{equation}
Here $u$ denotes the idele class of its original local embedding,
$a$ is the differential idele relating the global additive character
to the specified local characters, and $d_v^{\times}u$ and the
primed distributions are exactly those of
\texttt{TraceformulaH1} and \texttt{princvalrhobeta}.
In the canonical rational character used in SWQ6 one has $a_v=1$
at every place, so the displayed differential term evaluates with
$\log|a|=\log1=0$; it has not been absorbed into another term.
The prime sum includes infinity and every finite prime.

The prime notation is essential. In the author's definition a
distribution $\varrho_{e_v}$ extends multiplicative measure across
zero with Fourier transform satisfying
$\widehat\varrho_{e_v}(1)=0$, and the local term is
\[
 \left\langle\varrho_{e_v},
   \lambda\longmapsto
       \frac{h((\lambda+1)^{-1})}{|\lambda+1|_v}\right\rangle.
\]
For compactly supported tests this is the stated original definition;
for the full space $B$, SWQ28 uses the extension on that test space
in the cited trace theorem. It is not an assertion that the ordinary
integral through $u=1$ converges. No replacement of $f$ by
$|\cdot|^{1/2}f$, and no completed zeta function, is introduced
into SWQ23--28.

\paragraph{Exact scope and analytic limits.}
The closure in SWQ15 is the declared weighted Bruhat closure, not
an unspecified Hilbert, distributional, or pointwise closure.
The quotient, ideal, and sharp statements above are proved in this
domain. The scalar trace-class realization and its regularized
explicit formula retain the analytic hypotheses of the original
source theorems. SWQ23 does not supply a new proof of those input
theorems or a positivity result. It supplies a fully specified
same-label receiving object and proves how its operations and every
trace term descend. Its numerical values are exactly the original
ones: every such value is attained on top-labelled inputs, and every
labelled input has the displayed numerical amplitude. The support
data remain in $G_L(Q)$ and in $G_L(\mathbb C)$; an additive
numerical observation necessarily forgets them by SWQ20. The Bourne
quotient performs that forgetting already in its source, whereas
SWQ17 retains $e\ne\tau$ for every nontrivial $L$.
